%% kapa: exam=1 round=29 year=2018 period=1 form=A subject=civil_law subject_ko=민법 seq=11 number=11
%% source: Q-Net 감정평가사 기출문제 (cst00309) artlSeq=5202808 file=extracted/1008/5202808/2018년 제29회 감정평가사 1차시험 기출문제/2018년 제29회 감정평가사 1차시험 기출문제/1교시 A형.pdf
%% conversion: kordoc 4.24.51 -> tools/gen1.py (auto); review=proofread-sonnet; flags=none
\begin{question}{11}
대리권의 범위에 관한 설명으로 옳지 않은 것은? (다툼이 있으면 판례에 따름)
\begin{choices}
  \item 법정대리권의 범위는 법정대리인에 관한 규정에 의하여 결정된다.
  \item 임의대리권은 통상 그 권한에 부수하여 필요한 한도에서 상대방의 의사표시를 수령하는 대리권을 포함한다.
  \item 계약체결의 대리권을 수여받은 대리인은 특별한 사정이 없는 한 체결된 계약을 해제할 수 있는 권한을 갖지 않는다.
  \item 대리권의 범위를 정하지 않은 임의대리인은 대리의 목적인 물건의 성질이 변하지 않는 범위에서 그 이용행위를 할 수 있다.
  \item 예금계약의 체결을 위임받은 자의 대리권에는 특별한 사정이 없는 한 그 예금을 담보로 대출을 받을 수 있는 권한이 포함되어 있다.
\end{choices}
\end{question}
\answer{5}
